% !TeX root = ../papers/nc-p3.tex
\section{二つの二重点における割線の一階平方関係}
\label{sec:finite-two-double-five-space}
\begin{theorem}\label{thm:finite-two-double-five-space}
四辺が有限で全中央支持4の標準有限Frobenius窓において、点スキームの幾何学的支持が
十八点以上なら、反対階数10で少なくとも片側対称性があり、反対階数11は生じない。
\end{theorem}
\begin{proof}[合流する二節点の平方]
長さ20のGorenstein環なので長さ3以下の局所因子は曲線状である（接次元2で長さ3ならsocleが二次元）。
前節までの結果から十六単純点と二つの二重点 $d,e$ を扱えば十分である。
単純印 $p$ の五空間で、対称性がまだないとする。階数一切断が有限非被約なら
前節の十一割線の場合だけが残り、両超平面が全節点を避け、残余 $Z$ の別の支持も加えられない。
よって $J$ の外にある十一点のうち二つが二重点で、双方の一次方向が使われる。
そこにある割線は、非被約節点と単純節点を結ぶ異なる二本に限る。
実際、単純節点同士の割線は動かず、二重点そのものを切る接線も前節の有限平坦次数2の議論で動かない。
固定割線と超平面の交点は一意で、点座標の閉埋込みでは二重点を受け取れない。

四つの単純節点のベクトルを $u_1,\ldots,u_4$、二重節点を $v+\epsilon v'$ とする。
六つの列にただ一つある従属関係において $v'$ と四つの $u_i$ の係数は全て非零である。
これは定理\ref{thm:finite-marked-secant-pencil}の完全跡超平面が各局所socle上で非零であることによる。
代表の調整と $v'$ への $v$ の定数倍の加算により $v'=-\sum_i u_i$ とでき、
$v,u_1,\ldots,u_4$ は独立である。
$E_A$ の式 $a$ を $a(v)=1$ と規格化し $a_i=a(u_i)\ne0$,
$S_a=\sum a_i$, $x_i=v-u_i/a_i$ と置く。$x_i$ は $E_A$ の基底である。
四本の混合割線での超平面交点は $x_i$、六本の単純節点同士の割線では $x_i-x_j$、
二重節点での接線では $W_a=\sum a_i x_i$ となる。
二重点の一次方向を使う混合割線では、交点が非零の速度で
$x_i+\epsilon W_a$ の方向に動く。実際、$v+\epsilon v'$ を $a$ の値で割った微分は
$v'-a(v')v=W_a$ である。$b$ でも同じ表示をダッシュ付きの基底で使う。

十本の平方 $X_i=x_i^2$, $X_{ij}=(x_i-x_j)^2$ は対称二乗の基底である。
定理\ref{thm:finite-square-chart}の平方空間の単元倍での一致を双対化すると、
この十本を対応する十本の非零倍に送る線型同型が得られる。倍率を $s_i,s_{ij}$ とする。
この同型は十一番目の平方 $W_a^2$ を $cW_b^2$ に送り、微分した二本 $2x_i W_a$ を
それぞれ $2c_i x'_iW_b$ に $X'_i$ の倍を足したものに送る。
後者の不定項は各二重点での単元による代表の変更であり、$c,c_i\ne0$。
$W_a^2$ の $X_{jk}$ 係数は $-a_ja_k$ なので
\[
                   s_{jk}=c\frac{b_j b_k}{a_j a_k}.
\]
微分した式の $X_{ik}$ 係数は $-a_k$ ($k\ne i$) だから $c_i=c b_i/a_i$。
同じ式の $X_k$ 係数は $a_k$ ($k\ne i$) であり、$s_k a_k=c_i b_k$ となる。
微分する二添字を $1,2$ とすれば $k=3$ で比較して $b_1/a_1=b_2/a_2$。
従って使用された単純節点間の割線 $12$ において $a,b$ の交点が一致し、
そのグラフ点が $J$ の外にあることに反する。$S_a=0$ の場合も同じである。
\end{proof}
\begin{proof}[有限六節点と埋込み接点の印の取り直し]
切断が被約六点なら、命題\ref{prop:curvilinear-marked-pencils}の十三割線の議論で
両超平面は全節点を避ける。ここでは非単純点の\emph{全て}で二行列が剰余で比例する。
実際、剰余で独立なら局所環全域での鉛筆が一定割線となり、$B$ の値も超平面との交点に
固定される。これは非単純点を閉埋込みで写せない。この議論は任意の曲線状グラフにも成立する。

切断が正次元なら前節の証明により、未処理の場合は平面二次曲線の一つの外向き接点と
一つの孤立点で、$D$ 内の二支持、孤立点を通る六支持、外向き接線の三支持が各上界を使い切る。
$Z$ に別の支持が入らないことも同じである。孤立点の単根を通る相対鉛筆は全長で
その点を通り、外向き接線を持つ相対鉛筆も全長で固定接平面にあることは
前節の局所環 $k[[x,z]]/(xz,z^2)$ の計算による。従ってその六点と三点に二重点は入れられない。
二つの二重点はともに $D$ 内の鉛筆を与える二支持である。
この二つの剰余鉛筆の左核は同じ方向なので、核側の第一の位置で
$p,d,e$ は共線である。$d,e$ を先に固定すれば、この直線に入る単純印は
系\ref{cor:finite-no-four-on-line}により高々一つである。

階数10または11なら各単純印で少なくとも一方の和が五次元以下である。
共通四空間なら既に対称性が得られる。対称性がないとして、十六印のそれぞれについて
五空間の側を一つ選ぶ。埋込み接点のために除く印は側ごとに高々一つであり、
少なくとも十四印が有限六節点を与える。その七印以上で同じ側が選ばれる。
その側で $d$ を印に取り直すと七印の行列の剰余値が両四空間の交叉にある。
さらに $d$ 自身の二つの値は同じであり、それぞれの座標が閉埋込みなので
合計八つの異なる点が交叉に入る。どの平面切断も長さ7以下なので両四空間は一致する。
曲線状印での共通切断判定が対称性を与える矛盾である。対称二乗の次元10により階数11も排除される。
\end{proof}
